% note_draft_v1.tex — LaTeX conversion of the supervisor's draft v0.1
% (2026-07-22). Faithful to the markdown original (note_draft_v1.md);
% [Fable: ...] self-notes preserved as comments; audit-flagged passages
% carry "% [AUDIT n]" markers referring to AUDIT_v1.md — the prose itself
% is unmodified pending the supervisor's v2 decisions.
% Table in §3.1 is machine-generated by make_table.py (do not hand-edit).
\documentclass[11pt]{amsart}
\usepackage{amsmath,amssymb}
\usepackage[hidelinks]{hyperref}
\usepackage{booktabs}

\newtheorem{theorem}{Theorem}[section]
\newtheorem{lemma}[theorem]{Lemma}
\newtheorem{conjecture}[theorem]{Conjecture}
\theoremstyle{remark}
\newtheorem{remark}[theorem]{Remark}

\newcommand{\exc}{\operatorname{exc}}

\title[Measuring the moat]{Measuring the moat: the minimum excess of
Seymour orientations with large minimum out-degree}

% [AUDIT E: Cite-as name confirmed as Daiki Tahara — update placeholder]
\author{Daiki Tahara}
\address{Independent researcher, Osaka, Japan}
\thanks{In collaboration with Claude (Anthropic): Fable (supervision,
verification, literature) and Claude Code (implementation, search,
computation). Draft v0.1 --- 2026-07-22 --- for internal review
(Raa / Fable / Code).}

\begin{document}

\begin{abstract}
Seymour's second neighbourhood conjecture (SSNC) asserts that every oriented
graph has a vertex whose second out-neighbourhood is at least as large as its
first. A counterexample would be an oriented graph in which \emph{every}
vertex falls short. We study how close one can get. For an oriented graph $G$
define the \textbf{excess} $\exc(G) = \sum_v \max(0,\, d_2(v) - d_1(v) + 1)$,
where $d_i(v) = |N_i^+(v)|$; thus $\exc(G) = 0$ iff $G$ is a counterexample
to SSNC, and SSNC is equivalent to $\exc(G) \ge 1$ for all $G$. Since excess 1
is trivially attained using sinks, we impose the minimum out-degree bound
$\delta^+ \ge 8$ forced on counterexamples by Kaneko--Locke and the recent
$\delta^+=7$ elimination, and define $E(n) = \min\{\exc(G)\}$ over strongly
connected oriented graphs on $n$ vertices with $\delta^+ \ge 8$.

We give constructive upper bounds ($E(n) \le m$ whenever $n = mt$, $m \ge 3$,
$t \ge 8$, with a generalization to cycle-power skeletons covering many
further $n$), and a transplant construction derived from a stochastic-search
minimizer that yields $E(n) \le 5$ at several primes ($n = 47, 53, 59$). We
report machine-verified measurements for
% [AUDIT A4: witnesses exist for 25 <= n <= 75; no n=24 artifact]
$24 \le n \le 75$, all witnesses published with hashes and re-verifiable by
three independent implementations. On the lower-bound side we prove
$E(17) \ge 3$ by constraint-solver infeasibility, and report that the
measured floor at $n=50$ is invariant under raising the degree bound to
$\delta^+ \ge 9, 10$. Structural analysis of minimizers reveals two distinct
local architectures (a distributed ``ring membrane'' and a concentrated
``keystone cluster'') which nevertheless share identical local dynamics:
strict single-flip local minimality, large neutral plateaus, and stability
under all $\sim 1.7 \times 10^5$ coordinated two-flip attacks tested. We
conjecture a uniform lower bound $E(n) \ge 3$ and pose the regularity of
$n \mapsto E(n)$ as an open problem. All claims are reproducible from the
accompanying repository.
\end{abstract}

\maketitle

\section{Introduction}

Seymour's second neighbourhood conjecture \cite{DeanLatka1995} states:

\begin{conjecture}[SSNC]
Every oriented graph (digraph without 2-cycles) contains a vertex $v$ with
$|N_2^+(v)| \ge |N_1^+(v)|$.
\end{conjecture}

The conjecture is proved for tournaments \cite{Fisher1996,HavetThomasse2000}
and in many special classes; every vertex satisfies
$|N_2^+(v)| \ge \gamma\,|N_1^+(v)|$ for some vertex with
$\gamma = 0.715538\ldots$ \cite{HuangPeng2024}, improving
\cite{ChenShenYuster2003}. A counterexample must have minimum out-degree
$\delta^+ \ge 8$: out-degree $\le 6$ was eliminated by \cite{KanekoLocke2001}
and $\delta^+ = 7$ by \cite{SadhukhanSandeepSen2026} via a
signature-compression argument checked by CP-SAT.

Recent work has begun to study the \emph{equality boundary} of the
conjecture. \cite{Halkiewicz2026} calls a strongly connected oriented graph
\textbf{Pisa} if $\max_v (|N_2^+(v)| - |N_1^+(v)|) = 0$, i.e.\ no vertex
exceeds and at least one vertex meets the bound;
\cite{GuoKangZwaneveld2026} study \textbf{Seymour-tight} orientations, in
which \emph{every} vertex meets the bound, and develop a product calculus and
an algebraic classification for them. In the tight case the set of boundary
vertices is everything; in a generic Pisa graph it can be much smaller. This
note measures how small.

\subsection*{Definitions}
For an oriented graph $G$ write $d_1(v) = |N_1^+(v)|$, $d_2(v) = |N_2^+(v)|$,
and call $m(v) = d_2(v) - d_1(v)$ the \emph{margin} of $v$. Call $v$
\textbf{tight} if $m(v) = 0$ and \textbf{satisfied} if $m(v) \ge 0$ (a
``Seymour vertex''). Define
\[
\exc(G) \;=\; \sum_{v \in V(G)} \max\bigl(0,\; m(v) + 1\bigr).
\]
Then $\exc(G) = 0$ iff $G$ is a counterexample to SSNC; SSNC is the statement
$\exc \ge 1$; and for a Pisa graph the excess equals the number of tight
vertices. Adding a universal sink produces $\exc = 1$ trivially
\cite[\S1]{GuoKangZwaneveld2026}, so the interesting quantity requires the
degree bound that counterexamples must satisfy:
\[
E(n) \;=\; \min\{\, \exc(G) : |V(G)| = n,\ G \text{ strongly connected
oriented},\ \delta^+(G) \ge 8 \,\}.
\]
(Strong connectivity is justified by the standard minimal-counterexample
reduction: a closed out-section of a counterexample is a smaller
counterexample; cf.\ \cite{EspunyDiazGiraoGranetKronenberg2024}.)

$E(n) = 0$ for some $n$ iff SSNC is false. If SSNC is true, then
$E(n) \ge 1$ for all $n$, and any stronger uniform lower bound is a
strengthening of the conjecture. The contributions of this note:

\begin{enumerate}
\item \textbf{Upper bounds by construction} (\S\ref{sec:constructions}): a
  single ``out-set inclusion'' lemma yields both a family of
  \emph{pure-ring} constructions ($E(mt) \le m$ for $m \ge 3$, $t \ge 8$;
  generalized to cycle-power skeletons) and, after transplant of a
  search-discovered local structure, the bound $E(n) \le 5$ at the primes
  $n = 47, 53, 59$.
\item \textbf{Measurements} (\S\ref{sec:measurements}): a table of best
  known upper bounds for $24 \le n \le 75$
  % [AUDIT A4]
  with published, hash-identified witnesses; and a measurement-theoretic
  caveat we learned the hard way --- \emph{search attainment is not the
  floor} --- documented via a controlled experiment.
\item \textbf{Lower bounds and rigidity} (\S\ref{sec:lower}): $E(17) \ge 3$
  (CP-SAT infeasibility); invariance of the measured $n = 50$ floor under
  $\delta^+ \ge 9, 10$; and twelve one-hour UNKNOWNs for the
  tournament-restricted question $\exc \le 2$, $17 \le n \le 28$.
\item \textbf{Structure of minimizers} (\S\ref{sec:structure}): two
  architectures, one dynamics.
\item \textbf{Open problems} (\S\ref{sec:open}), foremost: does a uniform
  constant lower bound $E(n) \ge c \ge 2$ hold, and what governs the
  apparent floor $3 \le E(n) \le 5$?
\end{enumerate}

\subsection*{Verification protocol}
Every witness in this note is stored in the accompanying repository as an
adjacency list with a SHA-1 fingerprint, and every numerical claim is checked
by three implementations written independently (a numpy pipeline used by the
search; a pure-Python verifier; and a standard-library-only verifier written
from the definitions without access to the other two). A
\texttt{make verify-*} target reproduces each claim.

\section{Upper bounds by construction}\label{sec:constructions}

\subsection{The out-set inclusion lemma}

All constructions below are perturbed blow-ups of directed cycles. The
mechanism that keeps perturbations from leaking is the following observation.

\begin{lemma}[out-set inclusion]\label{lem:outset}
Let $G$ be an oriented graph and let $v, p$ be non-adjacent vertices with
$N_1^+(p) \subseteq N_1^+(v)$ and $p \notin N_1^+(v)$, $v \notin N_1^+(p)$.
Adding the arc $v \to p$ increases $d_1(v)$ by one and does not change
$d_2(v)$; consequently $m(v)$ decreases by exactly one. Moreover, for any
vertex $x$ with $v \in N_1^+(x)$ and $p \in N_1^+(x)$, the margins of $x$
are unchanged.
\end{lemma}

\begin{proof}
$N_2^+(v)$ gains at most the elements of
$N_1^+(p) \setminus (N_1^+(v) \cup \{v\}) = \emptyset$ and loses only $p$
itself if $p \in N_2^+(v)$; the hypothesis $N_1^+(p) \subseteq N_1^+(v)$
with $p$ previously at distance $\ge 2$\ldots
% [Fable: tighten this — state the exact bookkeeping for p in N_2^+(v) vs
% p at distance >= 3; in the layered applications below p in N_2^+(v)
% never occurs and the statement simplifies.]
% [AUDIT C9: if p in N_2^+(v), the margin drops by exactly 2, not 1 —
% d_1 gains 1 AND d_2 loses p. Case split is mandatory as stated.]
For observers $x$ as stated, $p$ already lies in $N_1^+(x)$, so no new
second neighbour is created.
\end{proof}

In words: \emph{if the target's out-set is already inside the shooter's
first neighbourhood, the shot costs the shooter one unit of margin and is
invisible to bystanders who see both.} The same arithmetic underlies the
in-star trick in the generalized lexicographic products of
\cite{GuoKangZwaneveld2026} --- and, as \S\ref{sec:structure} shows, it was
re-invented independently by stochastic search.

\subsection{Pure rings}

\begin{theorem}\label{thm:purering}
Let $m \ge 3$ and $t \ge 8$, and let $n = mt$. Then $E(n) \le m$.
\end{theorem}

\begin{proof}[Construction]
Partition $V$ into layers $L_0, \dots, L_{m-1}$ of size $t$; add all arcs
$L_i \to L_{i+1}$ (indices mod $m$); in each layer choose one \emph{pure}
vertex $p_i$ and add the arc $v \to p_i$ for every impure
$v \in L_i \setminus \{p_i\}$. Every impure vertex has
$d_1 = t + 1,\ d_2 = t$, hence margin $-1$; every pure vertex has
$d_1 = d_2 = t$, margin $0$. The graph is digon-free (arcs between layers
all point forward; inside a layer only $\to p_i$ arcs exist and $p_i$ sends
none), strongly connected, with $\delta^+ = t \ge 8$, and $\exc = m$.
\end{proof}

\begin{theorem}[cycle-power skeletons]\label{thm:power}
Let $m, k$ with $2k < m$, and $t \ge 1$ with $kt \ge 8$; let $n = mt$. Then
$E(n) \le m$.
\end{theorem}

\begin{proof}[Construction]
Replace the cycle skeleton by the $k$-th power $\vec{C}_m^{\,k}$: each layer
sends complete arcs to the next $k$ layers. Purity and the in-star are as
before; the out-set inclusion is now even cleaner, since the target's
out-set lies entirely inside the shooter's first neighbourhood. This covers
e.g.\ $n = 25\ (m{=}5,k{=}2,t{=}5)$, $n = 35\ (5,2,7)$, $n = 49\ (7,2,7)$.
[Verified numerically for $n = 25, 35, 49$; general proof identical to
Thm~\ref{thm:purering}.]
\end{proof}

Uneven layer sizes generally fail; the algebraic reason is the kernel
condition for generalized lexicographic products identified in
\cite[Thm 4.6]{GuoKangZwaneveld2026}, which for the plain cycle forces equal
sizes. Our constructions are \emph{not} Seymour-tight --- they sit at the
opposite extreme of the Pisa class, with as few tight vertices as the
skeleton allows.

\subsection{Keystone transplants: upper bounds at primes}\label{sec:keystone}

Ring constructions require $m \mid n$, and no admissible skeleton exists at
e.g.\ $n = 47, 53, 59$. Stochastic search (\S\ref{sec:measurements}) found a
minimizer at $n = 53$ with a different local architecture: a
\textbf{keystone cluster} --- a satisfied vertex $u$ of margin $+1$ together
with \emph{wing} vertices $w$ whose out-sets are exact copies of $N_1^+(u)$
plus the single arc $w \to u$ (verified equality of out-sets in the witness;
see \texttt{data/cluster\_structure\_n53.txt}). This is the out-set
inclusion mechanism again, asymmetric and ring-free.

Treating the cluster as an invariant module and deterministically deleting
(resp.\ cloning) bulk vertices yields witnesses establishing
\[
E(47) \le 5, \qquad E(53) \le 5, \qquad E(59) \le 5,
\]
each verified by the three-implementation protocol, and $\exc = 4$ at
$n = 57$ with \emph{no} $+1$ vertex (all four survivors tight) --- showing
the $+1$ ``insurance'' is not intrinsic to the family (at $n=57$ the ring
bound $E \le 3$ is anyway better).

\section{Measurements}\label{sec:measurements}

\subsection{Table of best known upper bounds}\label{sec:table}

The full table is machine-generated from the repository data
(\texttt{note/make\_table.py}); hashes are canonical-adjacency SHA-1
prefixes from the witness manifest.
% [AUDIT A3: witness identifiers below use graph hashes (manifest),
% resolving the code-lineage/graph-hash mixup of the draft's selected rows]

\begin{center}
\small
\textbf{Table 1.} Surveyed orders: the theorem against the constructions it supersedes.\par\smallskip
\begin{tabular}{rclll}
$n$ & $3+[3\nmid n]$ & earlier construction & fresh search & best stored witness \\ \hline
24 & 3 & 3 (ring m=3,k=1,t=8) & - & \texttt{2517130b} \\
25 & 4 & 5 (ring m=5,k=2,t=5) & 12 & \texttt{150003fd} \\
27 & 3 & 3 (ring m=3,k=1,t=9) & - & \texttt{5b5ae497} \\
30 & 3 & 3 (ring m=3,k=1,t=10) & 3 & \texttt{feba34e7} \\
35 & 4 & 5 (ring m=5,k=2,t=7) & 7 & \texttt{ca62560c} \\
40 & 4 & 4 (ring m=4,k=1,t=10) & 23 & \texttt{9cb88021} \\
45 & 3 & 3 (ring m=3,k=1,t=15) & 21 & \texttt{abd86a45} \\
47 & 4 & 5 (keystone transplant) & 9 & \texttt{ccc19f2e} \\
48 & 3 & 3 (ring m=3,k=1,t=16) & 13 (seed-free) & \texttt{d43ca1db} \\
49 & 4 & 7 (ring m=7,k=2,t=7) & 6 & \texttt{db2e4026} \\
50 & 4 & 5 (ring m=5,k=1,t=10) & 5 (20 seed-free) & \texttt{bce286aa} \\
51 & 3 & 3 (ring m=3,k=1,t=17) & - & \texttt{e5d2d5eb} \\
53 & 4 & 5 (keystone transplant) & 5 & \texttt{8562bd98} \\
57 & 3 & 3 (ring m=3,k=1,t=19) & - & \texttt{d02d73a9} \\
59 & 4 & 5 (keystone transplant) & 7 & \texttt{d90d86c3} \\
60 & 3 & 3 (ring m=3,k=1,t=20) & 33 & \texttt{44b4bca8} \\
75 & 3 & 3 (ring m=3,k=1,t=25) & 3 & \texttt{63346a8f} \\
\end{tabular}
\par\medskip
\textbf{Table 2.} The measured witnesses of \texttt{data/sonar\_best/}, summarised.\par\smallskip
\begin{tabular}{ll}
quantity & value \\ \hline
orders covered & 24..100, no gaps \\
stored witnesses & 77 (one per order) \\
excess = $3+[3\nmid n]$ & 77 of 77 orders \\
orders disagreeing & none \\
min out-degree observed & 8 (n=24) .. 33 (n=99) \\
strongly connected & 77 of 77 orders \\
\end{tabular}

\end{center}

\noindent Search values marked ``seed-free'' are the negative controls of
\S\ref{sec:notfloor}. At $n=57$ the transplant family also realises
$\exc=4$ (all survivors tight), which is not a bound record there.

\subsection{Search attainment is not the floor}\label{sec:notfloor}

We initially read fresh-search values at primes ($47{:}9$, $59{:}7$,
$53{:}5$) as a ``prime staircase''. A pre-registered controlled experiment
(fresh search on the \emph{composites} $n = 48, 50$ with all construction
seeds withheld) returned $13$ and $20$ --- far above the construction bounds
$3$ and $5$ at the same $n$. Conclusion: \textbf{fresh-search values measure
search reachability, not the floor}; the staircase was a measurement
artefact, later confirmed when the keystone transplant put all three primes
at $5$. We therefore report construction bounds and search attainments in
separate columns throughout. (Full logs, pre-registration commits, and the
negative-control seed runs are in the repository.)

\subsection{$\delta$-invariance}\label{sec:delta}

At $n = 50$ the measured floor $5$ is unchanged when the degree bound is
raised to $\delta^+ \ge 9$ and $\delta^+ \ge 10$ (the divisor-ring witness
has $\delta^+ = 10$ and satisfies all three constraints; forty generations
of search under each constraint produced no improvement). Any conjectured
lower bound may thus plausibly be sought in a form not depending on $\delta$
beyond the counterexample threshold.

\section{Lower bounds and rigidity}\label{sec:lower}

\begin{theorem}\label{thm:e17}
$E(17) \ge 3$.
\end{theorem}

\begin{proof}
CP-SAT infeasibility of the constraint system ``$n = 17$, oriented, strongly
connected, $\delta^+ \ge 8$, $\exc \le 2$'', with exact two-way reification
of the distance-two indicators; fixed seed; model and log published.
% [Code: resolved]
The model is \texttt{experiments/delta8/excess2\_search.py}
(\texttt{--n 17 --cap 2}, \texttt{random\_seed}~1, 8 workers), returning
INFEASIBLE in $218.9$~s; the machine-readable record is
\texttt{data/excess2\_results.jsonl}.
\end{proof}

For $18 \le n \le 22$ the same computation is UNKNOWN within 3600~s each;
the direct method's reach ends near $n = 17$.

\subsection*{Tournament rigidity}
By \cite{HavetThomasse2000}, sink-free tournaments have at least two
satisfied vertices; whether \emph{exactly two, both tight} is attainable in
a tournament (which would give $\exc = 2$ and refute any $E \ge 3$-type
bound in the tournament-dense regime) appears to be open. Our CP-SAT scans
of the tournament-restricted question $\exc \le 2$ for $17 \le n \le 28$
returned twelve consecutive UNKNOWNs at one hour each: the question is not
cheaply decidable in either direction. We record this as an explicit open
problem (\S\ref{sec:open}).
% [AUDIT D: the n=17 tournament instance is in fact decided (INFEASIBLE)
% by inheritance from Thm 4.1, tournaments being a subclass]

\section{Structure of minimizers: two architectures, one dynamics}
\label{sec:structure}

Two verified minimizers at $\exc = 5$:

\begin{itemize}
\item $n = 50$ (\texttt{17cd7dc8}): five tight survivors forming a
  \textbf{pentagon ring} $8 \to 16 \to 23 \to 36 \to 47 \to 8$ in which each
  survivor's first neighbourhood contains the next survivor and its second
  neighbourhood the one after; the sensitivity matrix (all valid single
  arc-flips) is diagonal --- killing any survivor promotes a reserve vertex
  from margin $-1$, an \emph{equal trade}.
  % [AUDIT A2: 4 of the 5 survivors admit an equal trade (dE=+0);
  % survivor v36's best flip costs +2 — "any" is too strong]
\item $n = 53$ (\texttt{dc922e89}): a \textbf{keystone cluster} --- margins
  $(+1, 0, 0, 0)$ on $\{18; 34, 35, 36\}$ with the wings' out-sets equal to
  the keystone's; the sensitivity matrix shows the only non-diagonal row at
  the keystone: killing it raises both wings.
\end{itemize}

Different architecture and different economy (the ring witness's
non-survivors all sit at margin exactly $-1$;
% [AUDIT A1: false — the n=50 witness's non-survivor margins are
% {-5:1, -4:1, -3:4, -2:17, -1:22}; only the tail depth differs (-5 vs -15)]
the keystone witness has a deep deficit tail down to $-15$) --- yet
identical dynamics: both are strict single-flip local minima with large
neutral plateaus, and the keystone's apparent $+1$ slack survives an
exhaustive two-flip attack ($171{,}288$ coordinated pairs around the
cluster; no decrease; $7{,}712$ neutral pairs). Whatever forbids
$\exc < 5$ at these $n$ is insensitive to the architecture, the economy,
and (by \S\ref{sec:delta}) the degree bound --- which suggests the true
obstruction is a global counting phenomenon rather than a local structural
one.

\begin{remark}[Pisa structure]
All minimizers above are Pisa graphs whose underlying undirected graphs are
irregular and far from $K_n$-minus-a-matching; together with small blow-ups
of directed cycles (e.g.\ the 2-blow-up of $\vec{C}_4$ on 8 vertices) they
show that the classification of underlying graphs of Pisa graphs is
substantially richer than conjectured in \cite[Conj.~5.1]{Halkiewicz2026}.
We have communicated explicit counterexamples to the author (private
communication, July 2026).
\end{remark}

\section{Open problems}\label{sec:open}

\begin{enumerate}
\item \textbf{Uniform floor.} Is there $c \ge 2$ with $E(n) \ge c$ for all
  $n$? All evidence above is consistent with $E(n) \ge 3$; we know only
  $E(17) \ge 3$ and, conditionally on SSNC, $E(n) \ge 1$.
\item \textbf{Regularity of $E$.} The best known upper bounds lie in
  $\{3,4,5,6\}$ for all measured $24 \le n \le 75$,
  % [AUDIT A4]
  with the value governed by available skeletons, not by primality. Is $E$
  eventually constant? Monotone in any sense?
\item \textbf{Tournaments with exactly two tight vertices.} Does a
  sink-free tournament with exactly two satisfied vertices, both tight,
  exist? (Twelve one-hour CP-SAT attempts, $17 \le n \le 28$, all UNKNOWN.)
\item \textbf{Keystone family limits.} The transplant family attains $4$
  (at $n=57$) with all survivors tight; is $3$ attainable ring-free? Is the
  family's native floor below the ring floor anywhere?
\item \textbf{Global counting.} Find a double-counting / matrix-analytic
  proof of any bound $E(n) \ge 2$. The matrix $S_D$ of
  \cite{GuoKangZwaneveld2026} expresses margins as row sums; excess is the
  positive part of $S_D \mathbf{1}$. The spectral theory of $S_D$ (untouched
  in the literature) and LP duality over flow relaxations both suggest
  themselves.
\end{enumerate}

\section{Methods and reproducibility}

Search: island-model program evolution (LLM mutation operators over
construction programs) and simulated annealing over a triangle/arc
representation;
% [AUDIT A5: no simulated annealing / triangle representation exists in
% this project's pipeline — island-model LLM evolution only]
all discovered improvements persisted to disk at discovery time with SHA-1
fingerprints. Verification: three independent implementations (\S1); every
claim maps to a \texttt{make verify-*} target. Pre-registration: search
predictions were committed to the repository before the corresponding runs
(commit hashes in \texttt{docs/}). Division of labour: direction and
decisions by the human author; supervision, literature, and independent
verification by Claude (Fable); implementation, search, and solver
engineering by Claude (Code). The mathematics should be judged on the
verifiability of the claims, not on the nature of the authors.

Repository: [URL on publication; private at draft time].

\bibliographystyle{amsalpha}
\bibliography{refs}

\end{document}
